\subsection{Идея стратегии с ограниченным обзором}

Зафиксируем размер алфавита $\sigma \geq 2$, целое $N \geq 1$ (\emph{размер окна}) и целое $\mathrm{MAX\_SH} \geq 0$ (\emph{максимальный сдвиг}). Пусть даны две независимые равномерно случайные строки $A, B$ над алфавитом $\Sigma_\sigma$ достаточно большой длины.

\emph{Стратегия с ограниченным обзором} --- это алгоритм, по шагам строящий общую подпоследовательность строк $A$ и $B$. На каждом шаге алгоритм имеет два указателя $i$ (по строке $A$) и $j$ (по строке $B$), при этом доступны:
\begin{itemize}
    \item окно следующих $N$ символов первой строки: $a_i, a_{i+1}, \ldots, a_{i+N-1}$;
    \item окно следующих $N$ символов второй строки: $b_j, b_{j+1}, \ldots, b_{j+N-1}$;
    \item текущий сдвиг $\mathrm{sh} \in [0, \mathrm{MAX\_SH}]$ (содержательно --- это нормированная разность $j - i$, ограниченная допустимым окном; см.~\S2.2 о нормировке).
\end{itemize}

Действия алгоритма на каждом шаге:
\begin{enumerate}
    \item Если $a_i = b_j$: алгоритм \emph{обязан} засчитать это совпадение: прибавить $+1$ к длине общей подпоследовательности и сдвинуть оба указателя на $1$, то есть $i \to i+1$, $j \to j+1$.
    \item Иначе $a_i \neq b_j$: алгоритм выбирает одно из двух действий:
    \begin{itemize}
        \item[(L)] сдвинуть только $i$: $i \to i+1$, $\mathrm{sh} \to \mathrm{sh}-1$;
        \item[(R)] сдвинуть только $j$: $j \to j+1$, $\mathrm{sh} \to \mathrm{sh}+1$.
    \end{itemize}
    Однако, если $\mathrm{sh}$ уже равен крайнему значению, разрешённое действие однозначно определено: например, при $\mathrm{sh} = \mathrm{MAX\_SH}$ запрещено действие (R).
\end{enumerate}

Длина общей подпоследовательности, построенной такой стратегией за $T$ шагов, обозначается $S_T^\pi$, где $\pi$ --- стратегия. По построению $S_T^\pi \leq \mathrm{lcs}(A_{1..T_A},\, B_{1..T_B})$, где $T_A, T_B$ --- финальные значения указателей.

\subsection{Кодирование состояния}

Каждый символ алфавита $\Sigma_\sigma$ представляется $\mathrm{BITS} = \lceil \log_2 \sigma \rceil$ битами. Окно из $N$ символов первой строки кодируется числом $a \in \{0, 1, \ldots, \sigma^N - 1\}$, в общем случае помещающимся в $\mathrm{MASK\_BITS} = \mathrm{BITS} \cdot N$ бит. Аналогично кодируется окно $b$ второй строки.

\emph{Нормировка сдвига.} В реализации сдвиг $\mathrm{sh}$ кодируется неотрицательным целым числом в диапазоне $[0, \mathrm{MAX\_SH}]$, представляющим текущую разность $j - i$ относительно базового опорного значения; параметр $\mathrm{MAX\_SH}$ задаёт допустимый размер этого диапазона. При таком соглашении границы $\mathrm{sh} = 0$ и $\mathrm{sh} = \mathrm{MAX\_SH}$ соответствуют выходу за допустимую полосу: при $\mathrm{sh} = 0$ запрещено действие (L), уменьшающее сдвиг, при $\mathrm{sh} = \mathrm{MAX\_SH}$ --- действие (R), увеличивающее сдвиг (см.~\S2.4). В коде это даёт массив ДП размера $\mathrm{MAX\_SH}+1$ слоёв.

\emph{Состояние ДП} --- тройка $s = (\mathrm{sh}, a, b)$. Общее число состояний составляет $(\mathrm{MAX\_SH}+1) \cdot \sigma^{2N}$.

\subsection{ДП-уравнения}

Обозначим через $f_K(\mathrm{sh}, a, b)$ ожидаемое значение длины общей подпоследовательности, построенной оптимальной стратегией за $K$ шагов, начиная с состояния $(\mathrm{sh}, a, b)$. Базовый случай:
\[
    f_0(\mathrm{sh}, a, b) = 0 \quad \text{для всех } (\mathrm{sh}, a, b).
\]

Для шага $K \geq 1$ рекуррентное соотношение имеет вид:
\[
    f_K(\mathrm{sh}, a, b) = \mathbb{E}_{a_0,\, b_0 \sim U(\Sigma_\sigma)}\bigl[ T_K(\mathrm{sh}, a, b, a_0, b_0) \bigr],
\]
где математическое ожидание берётся по новой паре символов $(a_0, b_0)$, поступающих в концы окон, а функция $T_K$ описывает оптимальный выбор действия:
\begin{align*}
T_K(\mathrm{sh}, a, b, a_0, b_0) =
\begin{cases}
1 + f_{K-1}(\mathrm{sh},\, a',\, b'),
    & \text{если } a \bmod \sigma = b \bmod \sigma, \\[2pt]
\max\bigl(f_{K-1}(\mathrm{sh}-1,\, a',\, b),\;
         f_{K-1}(\mathrm{sh}+1,\, a,\, b')\bigr),
    & \text{иначе,}
\end{cases}
\end{align*}
где $a'$ получается из $a$ удалением младшего символа и добавлением нового символа $a_0$ в старший разряд:
\[
    a' = (a \gg \mathrm{BITS}) \;|\; (a_0 \ll (\mathrm{MASK\_BITS} - \mathrm{BITS})),
\]
аналогично определяется $b'$.

В ветви совпадения (символы $a_i$ и $b_j$ совпадают, т.е. $a \bmod \sigma = b \bmod \sigma$) оба окна сдвигаются на одну позицию; в ветви несовпадения алгоритм выбирает, какое из окон сдвинуть.

\emph{Порядок принятия решения.} Подчёркиваем, что выбор действия в ветви несовпадения (то есть взятие $\max$) производится \emph{до} того, как становятся известны новые символы $a_0, b_0$: алгоритм видит лишь текущие окна $(a, b)$ и сдвиг $\mathrm{sh}$. Соответственно в формуле для $f_K$ математическое ожидание $\mathbb{E}_{a_0, b_0}[\,\cdot\,]$ берётся \emph{снаружи} от выражения $T_K$, содержащего $\max$. Так как $\max\bigl(f_{K-1}(\mathrm{sh}-1, a', b),\, f_{K-1}(\mathrm{sh}+1, a, b')\bigr)$ не зависит от $(a_0, b_0)$ в той ветви, где требуемое продолжение использует только одно из обновлённых окон (см.~уравнения выше), но в общем случае $a'$ и $b'$ являются функциями от $a_0$ и $b_0$, и порядок $\mathbb{E}[\max(\cdot)]$ (а не $\max(\mathbb{E}[\cdot])$) существенен.

\subsection{Граничные условия по сдвигу}

При $\mathrm{sh} = 0$ запрещено действие (L), уменьшающее сдвиг ниже нуля; при $\mathrm{sh} = \mathrm{MAX\_SH}$ запрещено действие (R), увеличивающее сдвиг сверх лимита. В этих случаях в несовпадающей ветви максимум вырождается:
\begin{itemize}
    \item при $\mathrm{sh} = 0$, $a \bmod \sigma \neq b \bmod \sigma$: $T_K = f_{K-1}(1,\, a,\, b')$ (доступно только действие (R));
    \item при $\mathrm{sh} = \mathrm{MAX\_SH}$, $a \bmod \sigma \neq b \bmod \sigma$: $T_K = f_{K-1}(\mathrm{MAX\_SH}-1,\, a',\, b)$ (доступно только действие (L)).
\end{itemize}

В реализации эти случаи учитываются заведением вспомогательных «зеркальных» массивов $m_0$ (для $\mathrm{sh} = 0$) и $m_1$ (для $\mathrm{sh} = \mathrm{MAX\_SH}$), хранящих значения предыдущего слоя в нужных «копиях».

\subsection{Универсализация на алфавит произвольного размера}

Для $\sigma$, не являющегося степенью двойки, не все значения битовой маски размера $\mathrm{MASK\_BITS}$ соответствуют допустимым окнам (например, при $\sigma = 3$ маска $\mathrm{0b11}$ соответствует символу $3 \notin \Sigma_3$). Чтобы избежать обхода недопустимых состояний, на этапе инициализации строится список \emph{валидных масок} $\mathcal{V} \subseteq \{0, 1, \ldots, 2^{\mathrm{MASK\_BITS}}-1\}$ размера $\sigma^N$.

При $\sigma = 9$, $N = 4$: $\mathrm{MASK\_BITS} = 16$, общее число масок $2^{16} = 65\,536$, из которых валидных --- $9^4 = 6\,561$.

\emph{Малый пример: $\sigma = 3$, $N = 5$.} Здесь $\mathrm{BITS} = 2$, $\mathrm{MASK\_BITS} = 10$, общее число масок $2^{10} = 1024$. Валидны те маски, в которых каждый из пяти двухбитовых блоков лежит в $\{00, 01, 10\}$; их $3^5 = 243$, то есть около $24\,\%$ от общего числа. Например, маска $\mathrm{0b}\,01\,10\,00\,01\,10$ (равная $390$) валидна и кодирует окно $(2, 1, 0, 2, 1)$, тогда как маска с блоком $11 = 3$ невалидна и в ДП-итерации не появляется.

\subsection{Длина симуляции}

При запуске стратегии симулируется некоторое число шагов $K$, которое играет роль «длины строки». Очевидно, $f_K(0, 0, 0) / K$ --- ожидаемая нормированная длина общей подпоследовательности оптимальной оконной стратегии при симуляции $K$ шагов. В разделе~3 будет показано, что предел этой величины при $K \to \infty$ существует и не превосходит $\gamma_\sigma$, что и обеспечивает строгость нижней оценки.

\subsection*{Выводы по разделу}
\addcontentsline{toc}{subsection}{Выводы по разделу 2}

\begin{enumerate}
    \item Сформулирована стратегия с ограниченным обзором: алгоритм видит окна из $N$ ближайших символов двух строк и текущий сдвиг, ограниченный значением $\mathrm{MAX\_SH}$.
    \item Состояние ДП $s = (\mathrm{sh}, a, b)$ имеет размер $(\mathrm{MAX\_SH}+1) \cdot \sigma^{2N}$.
    \item Рекуррентное соотношение явно выписано как математическое ожидание по двум независимым новым символам, поступающим в концы окон; граничные условия по сдвигу даны в явной форме.
    \item Универсализация на произвольный $\sigma$ достигается за счёт предвычисления списка валидных масок.
\end{enumerate}
